\documentclass[]{informeutn}

% Paquetes adicionales requeridos por la guía
\usepackage{csquotes}
\usepackage[backend=bibtex, style=alphabetic, sorting=ynt]{biblatex}
\addbibresource{fuentes.bib}
\usepackage{subcaption}
\usepackage{listings}
\usepackage{siunitx}

% Definición de unidades adicionales complementarias a informeutn.cls
\newcommand{\kHz}{\,\mathrm{kHz}}
\newcommand{\MHz}{\,\mathrm{MHz}}
\newcommand{\pF}{\,\mathrm{pF}}
\newcommand{\Vpp}{\,\mathrm{V}_{\mathrm{pp}}}
\newcommand{\mVpp}{\,\mathrm{mV}_{\mathrm{pp}}}
\newcommand{\uApp}{\,\mu\mathrm{A}_{\mathrm{pp}}}
\newcommand{\mApp}{\,\mathrm{mA}_{\mathrm{pp}}}

% Configuración del estilo de listados para simulaciones LTspice
\lstdefinestyle{ltspice}{
  backgroundcolor=\color{gray!10},
  basicstyle=\ttfamily\small,
  keywordstyle=\color{blue},
  frame=single,
  breaklines=true,
  captionpos=b
}

% Metadatos y datos de portada
\materia{Electrónica Aplicada I}
\titulo{Trabajo Práctico N°3}
\subtitulo{Amplificador en Configuración Emisor Común:\\
Diseño para Máxima Excursión Simétrica, Rectas de Carga y Pequeña Señal}
\comision{3R4}
\autores{Cortesini Pérez, Luciano Tomás -- 402719 \par
         Moreyra, Iván Agustín -- 413687 \par
         Pedregoza, Juan Ignacio -- 408644}
\fecha{Septiembre de 2026}

\begin{document}

  \maketitle
  \tableofcontents
  \setcounter{page}{1}
  \thispagestyle{plain}

  \chapter{Introducción y Objetivos}
\label{chap:introduccion}

La configuración en emisor común con transistor bipolar de juntura (BJT) constituye una de las topologías
amplificadoras fundamentales en la electrónica analógica discreta. Su amplia difusión se debe a su capacidad
de proveer simultáneamente una ganancia de tensión significativa y una ganancia de corriente apreciable,
brindando una ganancia de potencia global elevada con respecto a otras configuraciones elementales.

\section{Objetivos del Trabajo Práctico}

El presente trabajo de laboratorio persigue los siguientes objetivos técnicos y formativos:
\begin{enumerate}
  \item Diseñar un amplificador con transistor BJT NPN en configuración emisor común polarizado por divisor resistivo,
    estableciendo el punto de operación en reposo $Q(V_{CEQ}, I_{CQ})$ bajo el criterio de Máxima Excursión Simétrica (MES).
  \item Garantizar la estabilidad del punto de reposo frente a dispersiones del parámetro $\beta$ y variaciones
    de temperatura ambiental mediante la adecuada elección de la red de polarización y resistencia de emisor.
  \item Evaluar el comportamiento del circuito mediante simulación computarizada en LTspice, analizando el punto de
    operación en continua (.op) y la respuesta temporal en régimen transitorio (.tran) ante señales senoidales de $1\kHz$.
  \item Normalizar los componentes calculados adoptando resistores comerciales de las series estándar (E12),
    verificando que la desviación respecto a los valores analíticos no exceda la tolerancia máxima admitida del $10\%$.
  \item Realizar el análisis gráfico de las rectas de carga estática de corriente continua (CC) y dinámica de corriente
    alterna (CA), cuantificando la excursión simétrica real sin distorsión sobre la carga.
  \item Determinar mediante el modelo híbrido equivalente de pequeña señal los parámetros circuitales fundamentales:
    impedancia de entrada ($Z_i$), impedancia de salida ($Z_o$), ganancia de tensión ($A_v$) y ganancia de corriente ($A_i$).
  \item Montar físicamente el circuito en placa de circuito impreso (PCB) en laboratorio y relevar de manera experimental
    las magnitudes estáticas de polarización y los parámetros dinámicos empleando el método del resistor serie ($R_S$).
  \item Contrastar de forma cuantitativa los resultados analíticos, las simulaciones y las mediciones de laboratorio.
\end{enumerate}

\section{Especificaciones de Partida y Circuito Base}

Los parámetros iniciales del circuito fueron suministrados por la cátedra en la consigna de trabajo:
\begin{itemize}
  \item Tensión continua de alimentación: $V_{CC} = 15\V$.
  \item Resistencia de emisor para polarización: $R_E = 180\ohm$.
  \item Resistencia de colector: $R_C = 1.2\kohm$.
  \item Resistencia de carga conectada a la salida: $R_L = 1\kohm$.
  \item Frecuencia de señal de prueba de alterna: $f = 1\kHz$.
  \item Transistor NPN de silicio de uso general seleccionado: BC547B, con ganancia de corriente medida de $\beta = 284$.
  \item Caída de tensión directa base-emisor estimada para silicio: $V_{BE} = 0.70\V$.
\end{itemize}

En la figura~\ref{crkt:emisor_comun} se representa el diagrama esquemático completo de la etapa amplificadora a implementar.

\begin{figure}[!ht]
  \centering
  \begin{circuitikz}[scale=0.95, transform shape]
  % Referencia de masa
  \draw (0,0) node[ground] {}
    to[sV, l=$v_s$] (0,3)
    to[R, l=$R_s$] (1.8,3)
    to[C, l=$C_1$, -*] (3.5,3) coordinate (base_node);

  % Divisor resistivo de entrada
  \draw (3.5,6.5) node[vcc] {$V_{CC}$} coordinate (vcc_node)
    to[R, l=$R_2$, -*] (base_node)
    to[R, l=$R_1$, *-] (3.5,0) node[ground] {};

  % Transistor BJT NPN
  \draw (base_node) to[short] (4.2,3) node[npn, anchor=B] (Q1) {};
  \node[right=4pt] at (Q1.right) {$Q_1$ (BC546B)};

  % Malla de colector
  \draw (Q1.C) to[short, -*] ++(0,0.25) coordinate (coll_node)
    to[R, l_=$R_C$] ++(0,2) coordinate (aux) to [short] (vcc_node |- aux);

  % Malla de emisor con capacitor de desacoplo
  \draw (Q1.E) to[short, -*] ++(0,-0.5) coordinate (emit_node);
  \draw (emit_node) to[R, l=$R_E$] (5,0) node[ground] {};
\draw (emit_node) to[short] ++(1.5,0) coordinate(aux)
      to[C, l=$C_E$] (aux |- 0,0) node[ground] {};

  % Red de acoplo de salida y carga
  \draw (coll_node) to[C, l=$C_2$] ++(4,0) coordinate (out_node)
    to[short, -o] ++(1,0) node[above] {$v_L$};
  \draw (out_node) to[R, l=$R_L$, *-] (9,0) node[ground] {};
\end{circuitikz}

  \caption{Esquema circuital completo del amplificador en configuración emisor común con desacoplo de emisor.}
  \label{crkt:emisor_comun}
\end{figure}

El capacitor $C_1$ cumple la función de desacoplar el nivel continuo de la señal de entrada $v_s$, impidiendo que la
resistencia interna de la fuente modifique la polarización de base. De forma análoga, el capacitor de acoplo de salida
$C_2$ bloquea la tensión continua presente en el colector ($V_{CQ}$), asegurando que sobre la carga $R_L$ se desarrolle
únicamente la componente senoidal amplificada $v_L$. Finalmente, el capacitor de paso o desacoplo de emisor $C_E$ actúa
como un cortocircuito para la componente alterna a la frecuencia de trabajo ($1\kHz$), derivando la señal directamente a
masa y maximizando consecuentemente la ganancia de tensión de la etapa sin perjudicar la estabilidad en continua.

  \chapter{Diseño Analítico para MES}
\label{chap:diseno_teorico}

El objetivo prioritario del dimensionamiento en corriente continua consiste en situar el punto de trabajo en
reposo $Q(V_{CEQ}, I_{CQ})$ en el centro óptimo de la recta de carga dinámica de corriente alterna. De este modo se
garantiza que la señal senoidal de salida pueda oscilar con la máxima amplitud pico a pico antes de alcanzar la zona
de corte o la zona de saturación del transistor, condición denominada Máxima Excursión Simétrica (MES).

\section{Resistencias de Carga en Continua y Alterna}

Para el circuito bajo estudio, la resistencia equivalente que determina la recta de carga de corriente continua (CC)
comprende la suma de las resistencias conectadas en los terminales de colector y emisor:
\begin{equation}
  R_{CC} = R_C + R_E = 1200\ohm + 180\ohm = 1380\ohm
  \label{eq:Rcc}
\end{equation}

En régimen dinámico de corriente alterna (CA), a la frecuencia de interés ($1\kHz$), el capacitor de emisor $C_E$ actúa
como un cortocircuito perfecto a masa, anulando la realimentación negativa en alterna de $R_E$. A su vez, el capacitor
$C_2$ acopla la resistencia de carga $R_L$ en paralelo con $R_C$. Por consiguiente, la resistencia de carga dinámica
vista desde el colector resulta:
\begin{equation}
  R_{CA} = R_C \parallel R_L = \frac{R_C \cdot R_L}{R_C + R_L} = \frac{1200\ohm \cdot 1000\ohm}{1200\ohm + 1000\ohm} \approx 545.45\ohm
  \label{eq:Rca}
\end{equation}

\section{Condición de Máxima Excursión Simétrica}

La excursión máxima de la corriente de colector hacia la saturación viene dada por el valor de reposo:
\begin{equation}
  \Delta i_{C,\text{sat}} \approx I_{CQ} - I_{C,\text{sat}} \approx I_{CQ}
\end{equation}
asumiendo $V_{CE,\text{sat}} \approx 0\V$. Por su parte, la excursión hacia la región de corte está restringida por la
máxima caída de tensión admisible sobre la carga dinámica:
\begin{equation}
  \Delta i_{C,\text{corte}} = \frac{V_{CEQ}}{R_{CA}}
\end{equation}

Igualando ambas amplitudes para lograr una excursión perfectamente simétrica sin recorte:
\begin{equation}
  I_{CQ} = \frac{V_{CEQ}}{R_{CA}} \implies V_{CEQ} = I_{CQ} \cdot R_{CA}
  \label{eq:condicion_mes}
\end{equation}

Por otra parte, la ecuación de la malla de colector-emisor en continua impone:
\begin{equation}
  V_{CC} = V_{CEQ} + I_{CQ} \cdot R_{CC}
  \label{eq:malla_cc}
\end{equation}

Sustituyendo la condición de la ecuación~\eqref{eq:condicion_mes} en la ecuación~\eqref{eq:malla_cc}:
\begin{equation}
  V_{CC} = I_{CQ} \cdot R_{CA} + I_{CQ} \cdot R_{CC} = I_{CQ} \cdot (R_{CC} + R_{CA})
\end{equation}
de donde se despeja la corriente de reposo óptima para MES:
\begin{equation}
  I_{CQ(\text{MES})} = \frac{V_{CC}}{R_{CC} + R_{CA}} = \frac{15\V}{1380\ohm + 545.45\ohm} = \frac{15\V}{1925.45\ohm} \approx 7.790\mA
  \label{eq:Icq_mes}
\end{equation}

Reemplazando en la ecuación~\eqref{eq:condicion_mes}, la tensión continua colector-emisor resulta:
\begin{equation}
  V_{CEQ(\text{MES})} = I_{CQ(\text{MES})} \cdot R_{CA} = 7.790\mA \cdot 545.45\ohm \approx 4.250\V
  \label{eq:Vceq_mes}
\end{equation}

\section{Diseño de la Red de Polarización de Base}

Para asegurar la estabilidad del punto de trabajo ante la dispersión de la ganancia de corriente $\beta$ y las
variaciones de temperatura, se aplica el teorema de Thévenin al divisor resistivo de entrada ($R_1$ y $R_2$).
En la figura~\ref{crkt:thevenin} se muestra el circuito equivalente resultante.

\begin{figure}[!ht]
  \centering
  \begin{circuitikz}[scale=0.95, transform shape]
  % Malla Thévenin de base
  \draw (0,0) node[ground] {}
    to[battery1,invert, l=$V_{BB}$] (0,2.5)
    to[R, l=$R_B$, i=$I_{BQ}$] (2.8,2.5)
    to[short] (3.6,2.5) node[npn, anchor=B] (Q1) {};
  \node[right=4pt] at (Q1.right) {$Q_1$};

  % Colector
  \draw (Q1.C) to[R, l_=$R_C$, i_<=$I_{CQ}$] (Q1.C |- 0,5) node[vcc] {$V_{CC}$};

  % Emisor
  \draw (Q1.E) to[R, l=$R_E$, i=$I_{EQ}$] (Q1.E |- 0,0) node[ground] {};

  % Indicadores de tensión VBE y VCE
  \draw[dashed, color=gray] (3.6,2.4) to[open, v<=$V_{BE}$] (4.3,1.5);
  \draw[dashed, color=gray] (5.1,3.3) to[open, v^<=$V_{CEQ}$] (5,1.4);
\end{circuitikz}

  \caption{Circuito equivalente de Thévenin para el análisis de polarización de base.}
  \label{crkt:thevenin}
\end{figure}

El criterio de diseño estándar de la cátedra para garantizar que el divisor sea rígido y que la corriente que circula
por $R_1$ y $R_2$ sea mucho mayor que la corriente de base ($I_{BQ}$) establece:
\begin{equation}
  R_B \le \frac{\beta}{10} \cdot R_E
  \label{eq:criterio_rb}
\end{equation}
donde $\beta = 284$ es la ganancia medida del transistor BC547B. Adoptando el límite superior:
\begin{equation}
  R_B = \frac{284}{10} \cdot 180\ohm = 5112\ohm
  \label{eq:Rb_calculada}
\end{equation}

Analizando la malla de entrada del circuito equivalente de Thévenin:
\begin{equation}
  V_{BB} = I_{BQ} \cdot R_B + V_{BE} + I_{EQ} \cdot R_E
  \label{eq:malla_entrada}
\end{equation}

Teniendo en cuenta que $I_{BQ} = I_{CQ}/\beta$, el primer término puede reescribirse como:
\begin{equation}
  I_{BQ} \cdot R_B = \frac{I_{CQ}}{\beta} \cdot \left(\frac{\beta}{10} \cdot R_E\right) = 0.10 \cdot I_{CQ} \cdot R_E
\end{equation}

Aproximando $I_{EQ} \approx I_{CQ}$, la suma de las caídas resistivas resulta:
\begin{equation}
  I_{BQ} \cdot R_B + I_{EQ} \cdot R_E \approx 0.10 \cdot I_{CQ} \cdot R_E + 1.0 \cdot I_{CQ} \cdot R_E = 1.10 \cdot I_{CQ} \cdot R_E
\end{equation}
lo que conduce a la ecuación práctica de diseño:
\begin{equation}
  V_{BB} = I_{CQ(\text{MES})} \cdot 1.10 \cdot R_E + V_{BE}
  \label{eq:Vbb_formula}
\end{equation}

Evaluando numéricamente con $V_{BE} = 0.70\V$:
\begin{equation}
  V_{BB} = 7.790\mA \cdot 1.10 \cdot 180\ohm + 0.70\V = 1.542\V + 0.70\V = 2.242\V
  \label{eq:Vbb_calculada}
\end{equation}

\section{Determinación de los Resistores del Divisor}

Las ecuaciones que definen el divisor resistivo de entrada son:
\begin{equation}
  V_{BB} = V_{CC} \cdot \frac{R_1}{R_1 + R_2} \quad \text{y} \quad R_B = R_1 \parallel R_2 = \frac{R_1 \cdot R_2}{R_1 + R_2}
\end{equation}

Dividiendo miembro a miembro:
\begin{equation}
  \frac{V_{BB}}{R_B} = \frac{V_{CC}}{R_2} \implies R_2 = \left(\frac{V_{CC}}{V_{BB}}\right) \cdot R_B
  \label{eq:R2_formula}
\end{equation}

Reemplazando los valores obtenidos:
\begin{equation}
  R_2 = \left(\frac{15\V}{2.242\V}\right) \cdot 5112\ohm \approx 34.19\kohm
  \label{eq:R2_teorica}
\end{equation}

A partir de la relación del divisor, se obtiene la expresión para $R_1$:
\begin{equation}
  R_1 = \frac{R_2}{\left(\frac{V_{CC}}{V_{BB}}\right) - 1} = \frac{34.19\kohm}{6.690 - 1} \approx 6.01\kohm
  \label{eq:R1_teorica}
\end{equation}

En la tabla~\ref{tab:resumen_diseno_teorico} se sintetizan todas las magnitudes calculadas para el punto de reposo y los
componentes ideales del circuito bajo la condición de Máxima Excursión Simétrica.

\begin{table}[!ht]
  \centering
  \begin{tabular}{|l|c|c|l|}
    \hline
    \textbf{Parámetro} & \textbf{Símbolo} & \textbf{Valor Teórico} & \textbf{Unidad} \\
    \hline
    Resistencia Thévenin de base & $R_B$ & $5112$ & $\Omega$ \\
    Resistencia de carga en CC & $R_{CC}$ & $1380$ & $\Omega$ \\
    Resistencia de carga en CA & $R_{CA}$ & $545.45$ & $\Omega$ \\
    Corriente de colector de reposo (MES) & $I_{CQ(\text{MES})}$ & $7.790$ & $\mA$ \\
    Tensión colector-emisor de reposo & $V_{CEQ(\text{MES})}$ & $4.250$ & $\V$ \\
    Tensión Thévenin de base & $V_{BB}$ & $2.242$ & $\V$ \\
    Resistencia superior del divisor & $R_2$ & $34.19$ & $\kohm$ \\
    Resistencia inferior del divisor & $R_1$ & $6.01$ & $\kohm$ \\
    Tensión continua de emisor & $V_E$ & $1.402$ & $\V$ \\
    Tensión continua de base & $V_B$ & $2.102$ & $\V$ \\
    Tensión continua de colector & $V_C$ & $5.652$ & $\V$ \\
    Corriente de base de reposo & $I_{BQ}$ & $27.43$ & $\uA$ \\
    \hline
  \end{tabular}
  \caption{Resumen del dimensionamiento teórico para Máxima Excursión Simétrica.}
  \label{tab:resumen_diseno_teorico}
\end{table}

  \chapter{Normalización de Componentes y Recálculo del Punto Q}
\label{chap:normalizacion}

Los valores calculados analíticamente en el capítulo anterior ($R_1 = 6.01\kohm$ y $R_2 = 34.19\kohm$) corresponden
a valores teóricos ideales que no forman parte de las series comerciales normalizadas de resistores (series E12 y E24).
Por este motivo, se requiere adoptar componentes existentes en plaza asegurando que la dispersión introducida sobre
el punto de trabajo en reposo no exceda el límite del $10\%$ fijado en la consigna del trabajo práctico.

\section{Selección de Resistores Comerciales}

El factor primordial al seleccionar los resistores del divisor consiste en preservar con máxima fidelidad la relación
de transferencia de tensión Thévenin:
\begin{equation}
  \frac{V_{BB}}{V_{CC}} = \frac{R_1}{R_1 + R_2} \approx \frac{2.242\V}{15\V} \approx 0.1495
\end{equation}
manteniendo a su vez la resistencia equivalente $R_B = R_1 \parallel R_2$ en un valor cercano a los $5.11\kohm$ de diseño.

A partir de los valores comerciales disponibles se evaluaron dos alternativas:
\begin{itemize}
  \item Opción 1: $R_2 = 33\kohm$ y $R_1 = 5.6\kohm$, cuya relación resulta $5.6 / 38.6 \approx 0.1451$ (error del $-2.94\%$),
    con una resistencia Thévenin de $R_B = 4.78\kohm$.
  \item Opción 2: $R_2 = 39\kohm$ y $R_1 = 6.8\kohm$, cuya relación resulta $6.8 / 45.8 \approx 0.1485$ (error del $-0.67\%$),
    con una resistencia Thévenin de $R_B = 5.79\kohm$.
\end{itemize}

Se seleccionó la Opción 2 debido a su excelente aproximación a la relación de continua ideal y a que su resistencia de
base ($5.79\kohm$) respeta plenamente el criterio de estabilidad térmica frente a variaciones de $\beta$.

\section{Recálculo Analítico con Valores Normalizados}

Con los componentes comerciales adoptados ($R_1 = 6.8\kohm$, $R_2 = 39\kohm$, $R_C = 1.2\kohm$, $R_E = 180\ohm$ y
$R_L = 1\kohm$), se recalculan los parámetros de polarización en corriente continua.

El equivalente de Thévenin de la red de base pasa a ser:
\begin{equation}
  V_{BB} = 15\V \cdot \frac{6.8\kohm}{6.8\kohm + 39\kohm} = 15\V \cdot \frac{6.8}{45.8} \approx 2.227\V
  \label{eq:Vbb_norm}
\end{equation}
\begin{equation}
  R_B = \frac{6.8\kohm \cdot 39\kohm}{6.8\kohm + 39\kohm} \approx 5790.39\ohm
  \label{eq:Rb_norm}
\end{equation}

Planteando la ecuación exacta de la malla de entrada de base:
\begin{equation}
  I_{BQ} = \frac{V_{BB} - V_{BE}}{R_B + (\beta + 1) \cdot R_E} = \frac{2.227\V - 0.70\V}{5790.39\ohm + 285 \cdot 180\ohm} = \frac{1.527\V}{57090.39\ohm} \approx 26.75\uA
  \label{eq:Ibq_norm}
\end{equation}

A partir de la corriente de base se deduce la corriente de colector:
\begin{equation}
  I_{CQ} = \beta \cdot I_{BQ} = 284 \cdot 26.75\uA \approx 7.597\mA \approx 7.60\mA
  \label{eq:Icq_norm}
\end{equation}
y la corriente de emisor:
\begin{equation}
  I_{EQ} = (\beta + 1) \cdot I_{BQ} = 285 \cdot 26.75\uA \approx 7.624\mA
\end{equation}

Las tensiones en los diferentes nodos del circuito respecto a la referencia común de masa son:
\begin{equation}
  V_E = I_{EQ} \cdot R_E = 7.624\mA \cdot 180\ohm \approx 1.372\V
\end{equation}
\begin{equation}
  V_B = V_E + V_{BE} = 1.372\V + 0.70\V = 2.072\V
\end{equation}
\begin{equation}
  V_C = V_{CC} - I_{CQ} \cdot R_C = 15\V - (7.597\mA \cdot 1200\ohm) = 15\V - 9.116\V \approx 5.884\V
\end{equation}

Por ende, la tensión continua colector-emisor de reposo resulta:
\begin{equation}
  V_{CEQ} = V_C - V_E = 5.884\V - 1.372\V \approx 4.512\V
  \label{eq:Vceq_norm}
\end{equation}

Las corrientes a través de las ramas del divisor resistivo se calculan como:
\begin{equation}
  I_{R2} = \frac{V_{CC} - V_B}{R_2} = \frac{15\V - 2.072\V}{39\kohm} = \frac{12.928\V}{39\kohm} \approx 0.331\mA = 331.5\uA
\end{equation}
\begin{equation}
  I_{R1} = \frac{V_B}{R_1} = \frac{2.072\V}{6.8\kohm} \approx 0.305\mA = 304.7\uA
\end{equation}
verificando la ley de corrientes de Kirchhoff en el nodo de base:
\begin{equation}
  I_{R2} - I_{R1} = 331.5\uA - 304.7\uA = 26.8\uA \approx I_{BQ}
\end{equation}

\section{Verificación de Tolerancia}

Para certificar que los valores normalizados satisfacen las especificaciones de diseño, se calcula el error
porcentual relativo frente a los valores óptimos de MES:
\begin{equation}
  \varepsilon_{I_{CQ}} = \frac{|I_{CQ,\text{norm}} - I_{CQ,\text{MES}}|}{I_{CQ,\text{MES}}} \times 100\% = \frac{|7.60\mA - 7.79\mA|}{7.79\mA} \times 100\% = 2.44\%
\end{equation}
\begin{equation}
  \varepsilon_{V_{CEQ}} = \frac{|V_{CEQ,\text{norm}} - V_{CEQ,\text{MES}}|}{V_{CEQ,\text{MES}}} \times 100\% = \frac{|4.51\V - 4.25\V|}{4.25\V} \times 100\% = 6.12\%
\end{equation}

En la Cuadro~\ref{tab:comparativa_mes_normalizado} se expone la comparación directa entre el diseño MES y el circuito normalizado.

\begin{table}[!ht]
  \centering
  \resizebox{\textwidth}{!}{
  \begin{tabular}{|l|c|c|c|c|}
    \hline
    \textbf{Magnitud de Polarización} & \textbf{Diseño Teórico MES} & \textbf{Valor Normalizado} & \textbf{Error Relativo} & \textbf{Tolerancia} \\
    \hline
    Resistencia $R_1$ & $6.01\kohm$ & $6.80\kohm$ & $+13.1\%$ & Comercial E12 \\
    Resistencia $R_2$ & $34.19\kohm$ & $39.00\kohm$ & $+14.1\%$ & Comercial E12 \\
    Tensión Thévenin $V_{BB}$ & $2.242\V$ & $2.227\V$ & $-0.67\%$ & Admitida \\
    Resistencia Thévenin $R_B$ & $5.112\kohm$ & $5.790\kohm$ & $+13.3\%$ & Admitida \\
    Corriente de colector $I_{CQ}$ & $7.790\mA$ & $7.597\mA$ & $-2.44\%$ & $\le 10\%$ (Cumple) \\
    Tensión colector-emisor $V_{CEQ}$ & $4.250\V$ & $4.512\V$ & $+6.12\%$ & $\le 10\%$ (Cumple) \\
    Tensión de colector $V_C$ & $5.652\V$ & $5.884\V$ & $+4.10\%$ & $\le 10\%$ (Cumple) \\
    Tensión de emisor $V_E$ & $1.402\V$ & $1.372\V$ & $-2.14\%$ & $\le 10\%$ (Cumple) \\
    Tensión de base $V_B$ & $2.102\V$ & $2.072\V$ & $-1.43\%$ & $\le 10\%$ (Cumple) \\
    \hline
  \end{tabular}
  }
  \caption{Comparación cuantitativa entre el diseño ideal MES y el circuito con resistores normalizados.}
  \label{tab:comparativa_mes_normalizado}
\end{table}

Dado que tanto la corriente de colector como la tensión colector-emisor exhiben desviaciones significativamente
inferiores al límite exigido del $10\%$, se concluye que la etapa normalizada se encuentra en óptimas condiciones
para su simulación e implementación experimental.

  \input{chapters/5-simulacion.tex}
  \input{chapters/6-rectas_de_carga.tex}
  \chapter{Análisis de Pequeña Señal y Parámetros del Amplificador}
\label{chap:pequena_senal}

En este capítulo se desarrolla el modelo linealizado de pequeña señal para la configuración en emisor común a partir
de los parámetros híbridos del transistor bipolar, deduciendo analíticamente las expresiones de ganancia y de
impedancia. Asimismo, se describen los fundamentos de los métodos de medición experimental utilizados en laboratorio.

\section{Modelo Híbrido Equivalente del Transistor}

Para señales de amplitud reducida superpuestas al punto de reposo $Q$, el transistor BJT se comporta como un
dispositivo lineal. Despreciando los efectos de realimentación inversa ($h_{re} \approx 0$) y la conductancia de
salida ($h_{oe} \approx 0$) frente a la carga de colector, el transistor queda representado por su resistencia
dinámica de entrada $h_{ie}$ y su generador de corriente controlado por corriente $h_{fe} \cdot i_b$.

En la figura~\ref{crkt:pequena_senal} se presenta el circuito equivalente de pequeña señal del amplificador completo.

\begin{figure}[!ht]
  \centering
  \begin{circuitikz}[scale=0.9, transform shape]
  % Terminal de entrada
  \draw (0,0) node[ground] {}
    to[open, v^>=$v_i$] (0,2.5)
    to[short, o-*] (1.5,2.5) coordinate (in_node);

  % Resistencia RB = R1 // R2
  \draw (in_node) to[R, l=$R_B$, -] (1.5,0) node[ground] {};

  % Modelo híbrido del transistor
  \draw (in_node) to[short] (3.0,2.5)
    to[R, l=$h_{ie}$, i=$i_b$] (3.0,0) node[ground] {};

  % Fuente dependiente de corriente de colector
  \draw (5.5,0) node[ground] {}
    to[cisourceAM, l=$h_{fe} i_b$] (5.5,2.5)
    to[short, -] (7.0,2.5) coordinate (c_node);

  % Resistencia de colector RC
  \draw (7.0,2.5) to[R, l=$R_C$, -] (7.0,0) node[ground] {};

  % Resistencia de carga RL
  \draw (c_node) to[short] (8.5,2.5) coordinate (l_node)
    to[R, l=$R_L$, -] (8.5,0) node[ground] {};

   % 1. Dibujamos el terminal de salida
  \draw (l_node) to[short, -o] (9.5,2.5) node[above] {$v_L$} coordinate (out_node);

  % 2. Colocamos la masa en el punto inferior y medimos la tensión en abierto hacia arriba
  \draw (9.5,0) node[ground] {} to[open, v_>=$v_L$] (out_node);
\end{circuitikz}

  \caption{Modelo equivalente de pequeña señal en parámetros híbridos para el amplificador emisor común.}
  \label{crkt:pequena_senal}
\end{figure}

A partir de la física del semiconductor, la resistencia dinámica de la juntura base-emisor polarizada en directa es:
\begin{equation}
  r_e = \frac{V_T}{I_{CQ}}
  \label{eq:re_prima}
\end{equation}
donde $V_T = k T / q \approx 25.0\mV$ representa el potencial térmico a temperatura ambiente de laboratorio ($T \approx 290\Kelvin$).

La resistencia dinámica de entrada del transistor en emisor común ($h_{ie}$) resulta:
\begin{equation}
  h_{ie}  \approx h_{fe} \cdot \frac{V_T}{I_{CQ}} = 284 \cdot \frac{25.0\mV}{7.683\mA} \approx 924.08\ohm
  \label{eq:hie_calculo}
\end{equation}
y la ganancia de corriente en cortocircuito es $h_{fe} = \beta = 284$.

\section{Cálculo Analítico de Parámetros del Cuadripolo}

A partir del análisis de mallas y nodos del circuito de la figura~\ref{crkt:pequena_senal}, se deducen las expresiones
de los parámetros fundamentales del amplificador.

\subsection{Impedancia de Entrada (\texorpdfstring{$Z_i$}{Zi})}
La impedancia que la etapa presenta a la fuente de excitación comprende el paralelo de la red de polarización de base
($R_B = R_1 \parallel R_2$) y la resistencia de entrada del BJT ($h_{ie}$):
\begin{equation}
  Z_i = R_1 \parallel R_2 \parallel h_{ie} = R_B \parallel h_{ie}
  \label{eq:Zi_formula}
\end{equation}

Evaluando con $R_B = 5790.39\ohm$ y $h_{ie} = 924.08\ohm$:
\begin{equation}
  Z_i = \frac{5790.39\ohm \cdot 924.08\ohm}{5790.39\ohm + 924.08\ohm} \approx 796.90\ohm
  \label{eq:Zi_teorica}
\end{equation}

\subsection{Impedancia de Salida (\texorpdfstring{$Z_o$}{Zo})}
La impedancia de salida se evalúa anulando la fuente de excitación de entrada ($v_i = 0$). En dicha condición, la
corriente de base se extingue ($i_b = 0$), lo cual anula el generador dependiente de colector ($h_{fe} i_b = 0$).
Por lo tanto, la impedancia de salida vista desde los terminales de la carga resulta:
\begin{equation}
  Z_o = R_C = 1200\ohm
  \label{eq:Zo_teorica}
\end{equation}

\subsection{Ganancia de Tensión (\texorpdfstring{$A_v$}{Av})}
La tensión alterna desarrollada sobre la carga $R_L$ viene dada por la corriente inyectada por el colector sobre el
paralelo de resistencias:
\begin{equation}
  v_L = -h_{fe} \cdot i_b \cdot (R_C \parallel R_L)
\end{equation}
mientras que la tensión en los terminales de entrada es:
\begin{equation}
  v_i = i_b \cdot h_{ie}
\end{equation}

El cociente entre ambas tensiones define la ganancia de tensión de la etapa:
\begin{equation}
  A_v = \frac{v_L}{v_i} = -\frac{h_{fe} \cdot (R_C \parallel R_L)}{h_{ie}} = -\frac{R_C \parallel R_L}{r_e'}
  \label{eq:Av_formula}
\end{equation}

Reemplazando los valores numéricos:
\begin{equation}
  A_v = -\frac{545.45\ohm}{3.254\ohm} \approx -167.62
  \label{eq:Av_teorica}
\end{equation}
donde el signo negativo explicita la inversión de fase de $180^\circ$ de la topología emisor común.

\subsection{Ganancia de Corriente (\texorpdfstring{$A_i$}{Ai})}
La ganancia de corriente se define como la razón entre la corriente alterna que atraviesa la carga ($i_L$) y la
corriente total que suministra el generador a la entrada del amplificador ($i_i$):
\begin{equation}
  i_L = \frac{v_L}{R_L}, \quad i_i = \frac{v_i}{Z_i}
\end{equation}
\begin{equation}
  A_i = \frac{i_L}{i_i} = \frac{v_L / R_L}{v_i / Z_i} = A_v \cdot \frac{Z_i}{R_L}
  \label{eq:Ai_formula}
\end{equation}

En valor absoluto:
\begin{equation}
  |A_i| = 167.62 \cdot \frac{796.90\ohm}{1000\ohm} \approx 133.58
  \label{eq:Ai_teorica}
\end{equation}

\section{Métodos de Medición Experimental en Laboratorio}

La determinación práctica de las impedancias en pequeña señal no puede realizarse de forma directa con un ohmímetro
convencional, ya que los instrumentos de corriente continua alterarían los puntos de polarización del transistor. Por
ello, se recurre al método de inserción de una resistencia serie ($R_S$) calibrada.

\subsection{Medición de \texorpdfstring{$Z_i$, $A_v$ y $A_i$}{Zi, Av y Ai} (Método del Resistor Serie de Entrada)}
En la figura~\ref{crkt:medicion_zi} se ilustra el esquema de conexionado.

\begin{figure}[!ht]
  \centering
  \begin{circuitikz}[scale=0.9, transform shape]
  % Generador y resistencia serie
  \draw (0,0) node[ground] {}
    to[sV, l=$v_s$] (0,2.5)
    to[R, l=$R_S$, *-] (2.2,2.5) coordinate (vi_node)
    to[short, -o] (2.2,2.5);

  \draw (0,0) to[short] (2.2,0) ;
  \draw (0,2.5) to[short, *-o] (0,3.2) node[above] {CH1 ($v_s$)};
  \draw (vi_node) to[short, *-o] (2.2,3.2) node[above] {CH2 ($v_i$)};

  % Bloque amplificador
  \draw (2.7, -0.5) rectangle (5.5, 3.0);
  \node at (4.1, 1.25) [align=center] {Amplificador\\Emisor Común};

  \draw (vi_node) to[short] (2.7, 2.5);
  \draw (2.2, 0) to[short] (2.7, 0);

  % Carga y medición de salida
  \draw (5.5, 2.5) to[short, -*] (6.8, 2.5) coordinate (vl_node)
    to[R, l=$R_L$] (6.8, 0) node[ground] {};
  \draw (5.5, 0) to[short] (6.8, 0);
  \draw (vl_node) to[short, -o] (7.5, 2.5) node[right] {$v_L$ (Salida)};
\end{circuitikz}

  \caption{Esquema de medición de $Z_i$, $A_v$ y $A_i$ intercalando un resistor serie $R_S$ en la entrada.}
  \label{crkt:medicion_zi}
\end{figure}

Se intercala un resistor conocido $R_S = 820\ohm$ entre el generador de funciones y la entrada del amplificador.
Mediante un osciloscopio de dos canales se miden simultáneamente la tensión suministrada por el generador ($v_s$) y la
tensión efectiva que ingresa a la base tras la caída en $R_S$ ($v_i$).

La corriente que ingresa a la etapa resulta de la ley de Ohm sobre el resistor testigo:
\begin{equation}
  i_i = \frac{v_s - v_i}{R_S}
  \label{eq:ii_exp}
\end{equation}

Conociendo $i_i$ y $v_i$, la impedancia de entrada se determina de forma inmediata como:
\begin{equation}
  Z_i = \frac{v_i}{i_i} = \frac{v_i \cdot R_S}{v_s - v_i}
  \label{eq:Zi_exp_formula}
\end{equation}

Simultáneamente, midiendo la amplitud pico a pico de la tensión sobre la carga $R_L$ ($v_L$), se computan la ganancia
de tensión experimental:
\begin{equation}
  A_v = \frac{v_L}{v_i}
  \label{eq:Av_exp_formula}
\end{equation}
y la ganancia de corriente experimental:
\begin{equation}
  A_i = \frac{i_L}{i_i} = \frac{v_L / R_L}{(v_s - v_i) / R_S}
  \label{eq:Ai_exp_formula}
\end{equation}

\subsection{Medición de \texorpdfstring{$Z_o$}{Zo} (Método de Inyección en Salida)}
Para evaluar la impedancia de salida de la etapa, se utiliza el montaje de la figura~\ref{crkt:medicion_zo}.

\begin{figure}[!ht]
  \centering
  \begin{circuitikz}[scale=0.9, transform shape]

  % Bloque amplificador
  \draw (1.2, -0.5) rectangle (4.0, 3.0);
  \node at (2.6, 1.25) [align=center] {Amplificador\\Emisor Común};

  \draw (0.0, 2.0) to[short] (1.2, 2.0);
  \draw (0.0, 0) to[short] (1.2, 0);
  \draw (0.0, 0) to[short] (0, 2);
  \draw (0.6, 0) node[ground] {} ;

  % Red de inyección en salida
  \draw (4.0, 2.0) to[short, -*] (5.2, 2.0) coordinate (vo_node)
    to[R, l_=$R_S$] (7.2, 2.0)
    to[sV, l_=$v_s$] (7.2, 0) node[ground] {};
  \draw (4.0, 0) to[short] (7.2, 0);

  \draw (vo_node) to[short, *-o] (5.2, 3.0) node[above] {CH2 ($v_o$)};
  \draw (7.2, 2.0) to[short, *-o] (7.2, 3.0) node[above] {CH1 ($v_s$)};
\end{circuitikz}

  \caption{Esquema de inyección de señal para la medición de la impedancia de salida $Z_o$.}
  \label{crkt:medicion_zo}
\end{figure}

Se cortocircuita a masa la entrada de señal en alterna ($v_{in} = 0$) para anular cualquier excitación propia del
transistor, conservando intacta su polarización en continua. En el terminal de colector se desacopla la carga $R_L$ y se
conecta el generador de funciones en serie con un resistor patrón de $R_S = 1200\ohm$.

Midiendo la tensión del generador ($v_s$) y la tensión en el colector ($v_o$), la corriente inyectada es:
\begin{equation}
  i_o = \frac{v_s - v_o}{R_S}
\end{equation}
y la impedancia de salida queda determinada como:
\begin{equation}
  Z_o = \frac{v_o}{i_o} = \frac{v_o \cdot R_S}{v_s - v_o}
  \label{eq:Zo_exp_formula}
\end{equation}

  \chapter{Implementación Experimental y Mediciones de Laboratorio}
\label{chap:mediciones_laboratorio}

En este capítulo se expone el desarrollo experimental llevado a cabo en el laboratorio de Electrónica Aplicada I,
detallando el instrumental empleado, el procedimiento de medición tanto en corriente continua como en corriente
alterna, y el análisis comparativo consolidado entre las predicciones teóricas, las simulaciones y los valores medidos.

\section{Instrumental de Laboratorio Utilizado}

Para la caracterización del transistor y el relevamiento del circuito implementado sobre placa de circuito impreso (PCB)
se utilizó el siguiente equipamiento de laboratorio:
\begin{enumerate}
  \item Multímetro digital UNI-T UT890C: empleado para la verificación de componentes y la medición directa del
    parámetro estático de ganancia de corriente del transistor ($h_{FE} = \beta = 284$).
  \item Multímetro digital portátil Tenmars TM-86: empleado para el relevamiento de las tensiones nodales de continua en
    los terminales de emisor, base y colector.
  \item Generador de funciones de laboratorio FG-8002: utilizado como fuente de excitación senoidal a una frecuencia de
    $f = 1\kHz$ con baja distorsión armónica.
  \item Osciloscopio analógico de doble canal con calibración de ganancia: empleado para la medición de amplitudes pico
    a pico ($v_s, v_i, v_L, v_o$) y la verificación de la pureza espectral y ausencia de distorsión.
  \item Fuente de alimentación de corriente continua regulada: ajustada a una tensión de salida nominal de $15\V$
    (tensión real medida en bornes de alimentación: $V_{CC} = 14.77\V$).
\end{enumerate}

\section{Caracterización del Dispositivo y Montaje}

Previo al ensamblado del circuito, se identificaron los terminales del transistor BJT BC546B y se midió su ganancia de8
corriente en continua ($h_{FE}$) con el multímetro UNI-T UT890C en su escala específica, obteniéndose un valor de
$\beta = 284$ a temperatura ambiente de laboratorio, como se documenta en la figura~\ref{fig:foto_beta}.

\begin{figure}[!ht]
  \centering
  \includegraphics[width=0.45\textwidth]{pictures/Beta.jpeg}
  \caption{Medición del parámetro $\beta = 284$ del transistor BC5476 mediante el multímetro UNI-T UT890C.}
  \label{fig:foto_beta}
\end{figure}

El montaje se concretó sobre una placa de circuito impreso específica para el práctico de emisor común, disponiendo los
resistores comerciales ($R_1 = 6.8\kohm$, $R_2 = 39\kohm$, $R_C = 1.2\kohm$, $R_E = 180\ohm$ y $R_L = 1\kohm$) y los
capacitores electrolíticos de acoplo ($C_1 = C_2 = 22\uF$) y de desacoplo de emisor ($C_E = 2200\uF$).

\section{Mediciones en Corriente Continua (Polarización)}

Siguiendo las indicaciones metodológicas de la cátedra, todas las tensiones nodales fueron medidas con la punta negativa
del multímetro unida rígidamente a la referencia común de masa. Posteriormente, por leyes de Kirchhoff y Ohm, se
dedujeron las tensiones diferenciales y las corrientes de rama.

Los valores registrados con el multímetro Tenmars TM-86 fueron:
\begin{itemize}
  \item Tensión real de alimentación: $V_{CC} = 14.77\V$.
  \item Tensión de emisor respecto a masa: $V_E = 1.384\V$.
  \item Tensión de base respecto a masa: $V_B = 2.064\V$.
  \item Tensión de colector respecto a masa: $V_C = 5.550\V$.
\end{itemize}

A partir de estas lecturas se deducen:
\begin{equation}
  V_{BE} = V_B - V_E = 2.064\V - 1.384\V = 0.680\V
\end{equation}
\begin{equation}
  V_{CEQ} = V_C - V_E = 5.550\V - 1.384\V = 4.166\V
\end{equation}
\begin{equation}
  V_{RC} = V_{CC} - V_C = 14.77\V - 5.550\V = 9.220\V
\end{equation}
\begin{equation}
  V_{R2} = V_{CC} - V_B = 14.77\V - 2.064\V = 12.706\V
\end{equation}
\begin{equation}
  V_{R1} = V_B = 2.064\V
\end{equation}

Las corrientes correspondientes se obtienen aplicando la ley de Ohm sobre cada resistor:
\begin{equation}
  I_{CQ} = \frac{V_{RC}}{R_C} = \frac{9.220\V}{1200\ohm} \approx 7.683\mA
\end{equation}
\begin{equation}
  I_{EQ} = \frac{V_E}{R_E} = \frac{1.384\V}{180\ohm} \approx 7.689\mA
\end{equation}
\begin{equation}
  I_{R2} = \frac{V_{R2}}{R_2} = \frac{12.706\V}{39000\ohm} \approx 0.3258\mA = 325.8\uA
\end{equation}
\begin{equation}
  I_{R1} = \frac{V_{R1}}{R_1} = \frac{2.064\V}{6800\ohm} \approx 0.3035\mA = 303.5\uA
\end{equation}
\begin{equation}
  I_{BQ} = I_{R2} - I_{R1} = 325.8\uA - 303.5\uA = 22.27\uA
\end{equation}

En la figura~\ref{fig:fotos_dc} se muestran las lecturas de los instrumentos en la placa bajo ensayo, destacándose
la medición de $V_{CEQ} = 4.16\V$ y las lecturas de tensión de colector.

\begin{figure}[!ht]
  \centering
  \begin{subfigure}[t]{0.31\textwidth}
    \centering
    \includegraphics[width=\linewidth]{pictures/medicion_colector_emisor.jpeg}
    \caption{Medición directa de $V_{CEQ} = 4.16\V$.}
    \label{fig:foto_vce}
  \end{subfigure}
  \hfill
  \begin{subfigure}[t]{0.31\textwidth}
    \centering
    \includegraphics[width=\linewidth]{pictures/medicion_voltaje.jpeg}
    \caption{Medición de tensión sobre $R_C$.}
    \label{fig:foto_vrc}
  \end{subfigure}
  \hfill
  \begin{subfigure}[t]{0.31\textwidth}
    \centering
    \includegraphics[width=\linewidth]{pictures/medicion_voltaje2.jpeg}
    \caption{Medición de tensiones de reposo.}
    \label{fig:foto_voltaje2}
  \end{subfigure}
  \caption{Relevamiento de tensiones continuas de polarización en el circuito implementado en laboratorio.}
  \label{fig:fotos_dc}
\end{figure}

\section{Mediciones en Corriente Alterna (Pequeña Señal)}

\subsection{Medición de \texorpdfstring{$Z_i$, $A_v$ y $A_i$}{Zi, Av y Ai}}
Se conectó el generador de funciones ajustado a $1\kHz$ en serie con una resistencia testigo $R_S = 820\ohm$ conectada
a la entrada del amplificador. Con el atenuador del generador configurado en su mínima salida practicable para evitar
toda no linealidad, se observaron en el osciloscopio analógico las amplitudes pico a pico de las señales:
\begin{itemize}
  \item Tensión a la salida del generador (antes de $R_S$): $v_s = 18\mV_{pp}$.
  \item Tensión a la entrada de base del amplificador (después de $R_S$): $v_i = 10\mV_{pp}$.
  \item Tensión sobre la resistencia de carga $R_L$: $v_L = 1.30\V_{pp}$.
\end{itemize}

En la figura~\ref{fig:fotos_ac} se aprecian las capturas del osciloscopio analógico exhibiendo la señal senoidal
amplificada a la salida, con una escala de $0.5\V/\text{div}$ y ocupando aproximadamente $2.6$ divisiones verticales
($2.6 \times 0.5\V = 1.30\V_{pp}$).

\begin{figure}[!ht]
  \centering
  \begin{subfigure}[t]{0.48\textwidth}
    \centering
    \includegraphics[width=\linewidth]{pictures/medicion_ganacia_1.jpeg}
    \caption{Visualización de $v_L = 1.30\V_{pp}$ y hoja de cálculo.}
    \label{fig:foto_osc1}
  \end{subfigure}
  \hfill
  \begin{subfigure}[t]{0.48\textwidth}
    \centering
    \includegraphics[width=\linewidth]{pictures/medicion_ganancia_2.jpeg}
    \caption{Panel del osciloscopio y generador de funciones FG-8002.}
    \label{fig:foto_osc2}
  \end{subfigure}
  \caption{Verificación experimental de la señal de salida senoidal a $1\kHz$ sin distorsión visible.}
  \label{fig:fotos_ac}
\end{figure}

Aplicando las expresiones del método serie deducidas en la ecuación~\eqref{eq:ii_exp}:
\begin{equation}
  i_i = \frac{v_s - v_i}{R_S} = \frac{18\mV_{pp} - 10\mV_{pp}}{820\ohm} = \frac{8\mV_{pp}}{820\ohm} \approx 9.756\uA_{pp}
\end{equation}

A partir de esta corriente se calculan la impedancia de entrada y las ganancias experimentales:
\begin{equation}
  Z_i = \frac{v_i}{i_i} = \frac{10\mV_{pp}}{9.756\uA_{pp}} \approx 1025.0\ohm
\end{equation}
\begin{equation}
  A_v = \frac{v_L}{v_i} = \frac{1.30\V_{pp}}{10.0\mV_{pp}} = 130.0
\end{equation}
\begin{equation}
  A_i = \frac{i_L}{i_i} = \frac{v_L / R_L}{i_i} = \frac{1.30\V_{pp} / 1000\ohm}{9.756\uA_{pp}} = \frac{1.30\mA_{pp}}{0.009756\mA_{pp}} \approx 133.25
\end{equation}

\subsection{Medición de \texorpdfstring{$Z_o$}{Zo}}
Pasivando la entrada de señal ($v_{in} = 0$) y excitando la salida del amplificador mediante el generador conectado en
serie con un resistor $R_S = 1200\ohm$, se registraron las siguientes amplitudes pico a pico:
\begin{itemize}
  \item Tensión del generador: $v_s = 1.00\V_{pp}$.
  \item Tensión desarrollada en el colector: $v_o = 0.50\V_{pp}$.
\end{itemize}

La corriente inyectada sobre el colector fue:
\begin{equation}
  i_o = \frac{v_s - v_o}{R_S} = \frac{1.00\V_{pp} - 0.50\V_{pp}}{1200\ohm} = \frac{0.50\V_{pp}}{1200\ohm} \approx 0.4167\mA_{pp}
\end{equation}
resultando una impedancia de salida experimental idéntica a la teórica:
\begin{equation}
  Z_o = \frac{v_o}{i_o} = \frac{0.50\V_{pp}}{0.4167\mA_{pp}} = 1200\ohm
\end{equation}

\section{Tabla Comparativa Consolidada y Análisis de Errores}

En el Cuadro~\ref{tab:consolidada_general} se resumen y comparan los valores obtenidos a través de las cuatro etapas del
estudio: diseño analítico para MES, cálculo analítico con componentes normalizados, simulación computarizada en LTspice
y medición experimental en laboratorio.

\begin{table}[!ht]
  \centering
  \resizebox{\textwidth}{!}{
  \begin{tabular}{|l|c|c|c|c|c|}
    \hline
    \textbf{Magnitud Eléctrica} & \textbf{Teórico MES} & \textbf{Teórico Normalizado} & \textbf{Simulación LTspice} & \textbf{Medición Laboratorio} & \textbf{Error Exp. vs Norm.} \\
    \hline
    Tensión de alimentación $V_{CC}$ & $15.00\V$ & $15.00\V$ & $15.00\V$ & $14.77\V$ & $-1.53\%$ \\
    Tensión colector-emisor $V_{CEQ}$ & $4.250\V$ & $4.512\V$ & $4.610\V$ & $4.166\V$ & $-7.67\%$ \\
    Corriente de colector $I_{CQ}$ & $7.790\mA$ & $7.597\mA$ & $7.525\mA$ & $7.683\mA$ & $+1.13\%$ \\
    Tensión de emisor $V_E$ & $1.402\V$ & $1.372\V$ & $1.360\V$ & $1.384\V$ & $+0.87\%$ \\
    Tensión de base $V_B$ & $2.102\V$ & $2.072\V$ & $2.006\V$ & $2.064\V$ & $-0.39\%$ \\
    Tensión de colector $V_C$ & $5.652\V$ & $5.884\V$ & $5.970\V$ & $5.550\V$ & $-5.68\%$ \\
    Corriente de base $I_{BQ}$ & $27.43\uA$ & $26.75\uA$ & $30.90\uA$ & $22.27\uA$ & $-16.7\%$ \\
    Impedancia de entrada $Z_i$ & $796.9\ohm$ & $796.9\ohm$ & -- & $1025.0\ohm$ & $+28.6\%$ \\
    Impedancia de salida $Z_o$ & $1200\ohm$ & $1200\ohm$ & -- & $1200.0\ohm$ & $0.00\%$ \\
    Ganancia de tensión $|A_v|$ & $167.6$ & $167.6$ & $139.1$ & $130.0$ & $-22.4\%$ \\
    Ganancia de corriente $|A_i|$ & $133.6$ & $133.6$ & -- & $133.25$ & $-0.26\%$ \\
    \hline
  \end{tabular}
  }
  \caption{Comparativa consolidada de resultados teóricos, simulados y experimentales.}
  \label{tab:consolidada_general}
\end{table}

El análisis de la tabla~\ref{tab:consolidada_general} permite constatar un grado de concordancia notable en las magnitudes
de polarización en continua, donde todas las variables se mantuvieron dentro de márgenes de error inferiores al $8\%$.
En régimen dinámico, la ganancia de corriente experimental ($A_i = 133.25$) coincidió de forma casi exacta con el valor
teórico ($133.58$), con un error del $-0.26\%$, al igual que la impedancia de salida ($Z_o = 1200\ohm$, error del $0\%$).

Las discrepancias observadas en $Z_i$ y $A_v$ responden a factores físicos bien identificados:
\begin{itemize}
  \item La ganancia de tensión teórica ($167.6$) asume el modelo ideal $A_v = -R_{CA}/r_e'$, donde la resistencia
    intrínseca $r_e' \approx 3.25\ohm$ no contempla la resistencia de contacto óhmico de emisor y base del transistor
    físico ($r_{bb'}$ y $r_e$), las cuales reducen levemente la transconductancia efectiva ($g_m$). Al comparar el valor
    medido ($A_v = 130$) con la simulación que sí incorpora estos parámetros no ideales ($A_{v,\text{sim}} = 139.1$),
    el error relativo disminuye a un moderado $6.5\%$.
  \item En la medición de $Z_i$, la atenuación sobre la resistencia serie $R_S = 820\ohm$ involucra lecturas de muy
    pequeña amplitud en el osciloscopio analógico ($8\mV$ a $18\mV$), donde la apreciación visual de las divisiones de la
    pantalla y el ruido ambiental introducen una incertidumbre experimental estimada en $\pm 1\mV$, suficiente para
    justificar la diferencia hallada.
\end{itemize}

  \input{chapters/9-conclusiones.tex}

  % Bibliografía
  \nocite{*}
  \printbibliography[heading=bibnumbered, title={Fuentes}]

  % Anexos
  \chapter{Anexos}
\label{chap:anexos}

\section{Extracto de la Hoja de Datos del Transistor BC546B}

El transistor utilizado en los ensayos es un BJT NPN epitaxial de silicio de uso general, fabricado por ON
Semiconductor. En el Cuadro ~\ref{tab:datasheet_ratings} y la tabla~\ref{tab:datasheet_caract} se transcriben los valores
límites absolutos y las características eléctricas más relevantes extraídas de la hoja de datos oficial del fabricante.

\begin{table}[!ht]
  \centering
  \resizebox{\textwidth}{!}{
  \begin{tabular}{|l|c|c|c|}
    \hline
    \textbf{Parámetro Máximo Admisible} & \textbf{Símbolo} & \textbf{Valor Límite} & \textbf{Unidad} \\
    \hline
    Tensión colector-emisor con base abierta & $V_{CEO}$ & $45$ & $\V$ \\
    Tensión colector-base con emisor abierto & $V_{CBO}$ & $50$ & $\V$ \\
    Tensión emisor-base con colector abierto & $V_{EBO}$ & $6.0$ & $\V$ \\
    Corriente continua de colector máxima & $I_C$ & $100$ & $\mA$ \\
    Corriente pico de colector & $I_{CM}$ & $200$ & $\mA$ \\
    Disipación total de potencia a $T_A = 25^\circ\text{C}$ & $P_D$ & $625$ & $\mW$ \\
    Rango de temperatura de juntura operativa & $T_J$ & $-55 \text{ a } +150$ & $^\circ\text{C}$ \\
    \hline
  \end{tabular}
  }
  \caption{Límites máximos absolutos de operación del transistor BC546B (ON Semiconductor).}
  \label{tab:datasheet_ratings}
\end{table}

\begin{table}[!ht]
  \centering
  \resizebox{\textwidth}{!}{
  \begin{tabular}{|l|c|c|c|c|c|}
    \hline
    \textbf{Característica Eléctrica} & \textbf{Símbolo} & \textbf{Mín.} & \textbf{Típ.} & \textbf{Máx.} & \textbf{Unidad} \\
    \hline
    Ganancia de corriente en continua ($I_C = 2\mA, V_{CE} = 5\V$) & $h_{FE}$ & $200$ & $290$ & $450$ & -- \\
    Tensión de saturación colector-emisor ($I_C = 10\mA, I_B = 0.5\mA$) & $V_{CE(\text{sat})}$ & -- & $0.09$ & $0.25$ & $\V$ \\
    Tensión base-emisor en conducción ($I_C = 2\mA, V_{CE} = 5\V$) & $V_{BE(\text{on})}$ & $0.58$ & $0.66$ & $0.70$ & $\V$ \\
    Frecuencia de transición ($I_C = 10\mA, V_{CE} = 5\V, f = 100\MHz$) & $f_T$ & $100$ & $300$ & -- & $\MHz$ \\
    Capacitancia de colector a base ($V_{CB} = 10\V, f = 1\MHz$) & $C_{obo}$ & -- & $3.5$ & $4.5$ & $\text{pF}$ \\
    \hline
  \end{tabular}
  }
  \caption{Características eléctricas típicas del grupo de ganancia B del transistor BC546.}
  \label{tab:datasheet_caract}
\end{table}

La ganancia medida en el laboratorio ($\beta = 284$) se ubica en el centro de la distribución típica del grupo B ($290$).
Asimismo, la disipación estática de reposo en el circuito implementado resulta:
\begin{equation}
  P_{D,Q} = V_{CEQ} \cdot I_{CQ} = 4.166\V \cdot 7.683\mA \approx 32.01\mW
\end{equation}
valor que representa apenas el $5.1\%$ de la disipación máxima admisible ($625\mW$), garantizando que el componente
funciona holgadamente dentro de su zona de operación segura (SOA) y sin elevación perceptible de temperatura.

\section{Netlist Completo de Simulación en LTspice}

A continuación se reproduce el netlist de simulación correspondiente al circuito con valores comerciales normalizados:

\begin{lstlisting}[style=ltspice, caption={Netlist de simulación para el circuito con valores comerciales en LTspice.}, label={list:netlist_comercial}]
* C:\users\corte\Documents\utn\Electr-nica-Aplicada-\TP3\sim\ValoresComerciales.asc
Q1 n001 B n004 0 BC546C
VCC n003 0 15
Rc n003 n001 1.2k
Re n004 0 180
R2 n003 B 39k
R1 B 0 6.8k
C1 n004 0 2200uF
C2 n001 n005 22uF
Rl n005 0 1k
VBB n002 0 SINE(0 10m 1k)
C3 B n002 22uf
.model BC546C NPN(Is=21.43f Xti=3 Eg=1.11 Vaf=74.03 Bf=290
+ Ise=43.11f Ne=1.442 Ikf=97.09m Nk=.5 Xtb=1.5 Br=9.837
+ Isc=144.3f Nc=1.522 Ikr=2.476 Rc=1 Cjc=5.25p Mjc=.3147
+ Vjc=.5697 Fc=.5 Cje=11.5p Mje=.6115 Vje=.8194 Tr=10n
+ Tf=443.2p Itf=.7 Vtf=4 Xtf=4 Vceo=65 Icrating=100m mfg=Philips)
.tran 0 1 990m
.backanno
.end
\end{lstlisting}

\section{Planilla de Registro de Mediciones de Banco}

En la tabla~\ref{tab:planilla_cruda} se detallan los registros primarios volcados en la hoja de cálculo de laboratorio
\texttt{mediciones.ods}.

\begin{table}[!ht]
  \centering
  \resizebox{\textwidth}{!}{
  \begin{tabular}{|l|c|c|l|}
    \hline
    \textbf{Registro / Parámetro de Banco} & \textbf{Valor} & \textbf{Unidad} & \textbf{Observaciones de Laboratorio} \\
    \hline
    Tensión de alimentación medida $V_{CC}$ & $14.77$ & $\V$ & Medida en bornes de placa con multímetro \\
    Tensión de emisor $V_E$ & $1.384$ & $\V$ & Tensión continua respecto a masa \\
    Tensión de base $V_B$ & $2.064$ & $\V$ & Tensión continua respecto a masa \\
    Tensión de colector $V_C$ & $5.550$ & $\V$ & Tensión continua respecto a masa \\
    Caída sobre $R_C$ ($V_{RC}$) & $9.220$ & $\V$ & Calculada: $V_{CC} - V_C$ \\
    Tensión colector-emisor $V_{CEQ}$ & $4.166$ & $\V$ & Calculada: $V_C - V_E$ (multímetro: $4.16\V$) \\
    Corriente de colector $I_{CQ}$ & $7.683$ & $\mA$ & Deducida por ley de Ohm: $V_{RC} / R_C$ \\
    Corriente de base $I_{BQ}$ & $22.27$ & $\uA$ & Deducida: $I_{R2} - I_{R1}$ \\
    Ganancia medida en multímetro $\beta$ & $284$ & -- & Medida en zócalo hFE de multímetro UT890C \\
    Resistencia serie de entrada $R_S$ & $820$ & $\Omega$ & Resistor testigo para medición de $Z_i$ y $A_v$ \\
    Tensión del generador $v_s$ (entrada) & $18$ & $\mV_{pp}$ & Medida con osciloscopio analógico CH1 \\
    Tensión de entrada en base $v_i$ & $10$ & $\mV_{pp}$ & Medida con osciloscopio analógico CH2 \\
    Tensión de salida en carga $v_L$ & $1.30$ & $\V_{pp}$ & Señal senoidal limpia a $1\kHz$ \\
    Resistencia serie de salida $R_S$ & $1200$ & $\Omega$ & Resistor testigo para inyección en $Z_o$ \\
    Tensión del generador $v_s$ (salida) & $1.00$ & $\V_{pp}$ & Inyección en colector con entrada a masa \\
    Tensión de salida atenuada $v_o$ & $0.50$ & $\V_{pp}$ & Medida sobre terminal de colector \\
    \hline
  \end{tabular}
  }
  \caption{Planilla de datos crudos y observaciones experimentales registradas en banco de trabajo.}
  \label{tab:planilla_cruda}
\end{table}


\end{document}
